% a06_BALtorus_ct2: Step 3(E), second part: the block kernels at zero background (slice ladder, remainder, order-2
% infrared remainders), the multi-cut expansion, the cut-line mechanism, and dm^2_{K,k} = O((L^k eps)^{-1}) (Ba3 (1.28)).

\subsubsection*{The block kernels at zero background}

\emph{Setting} \cite[(1.25)--(1.28), p.~418]{Ba3}. On a torus $T_\eps$ aligned with the blocks (all sides $\ge1$), at $A=0$ and in $\eps$-units, put $l_j\defeq L^j\eps=t_j$, $G_j\defeq(C_0^{-1}+a_jl_j^{-2}P_j)^{-1}$ for $1\le j\le K$, with $P_j$ the orthogonal projection onto the functions constant on the blocks of side $l_j$ (so $G_K$ is the operator of \cref{lem:a06:P}) and $G_0\defeq0$; $S_j\defeq G_{j+1}-G_j$ (the slice of index $j$; $S_0=G_1$) and $R_j\defeq C_0-G_j$. Then $G_K=G_k+G_{K,k}$ with $G_{K,k}=\sum_{j=k}^{K-1}S_j$, which is \cite[(1.27)]{Ba3}, and $R_K$ is the ``operator with a regular kernel'' of \cite[(1.25)]{Ba3}: $R_K=G_K(a_Kl_K^{-2}P_K)C_0=C_0(a_Kl_K^{-2}P_K)G_K$. The vector kernels $G^v_j$, $S^v_j$, $R^v_j$ are the same with $\lmu$ in place of $m$, componentwise (Feynman gauge; \cite[p.~608]{Ba1}: at $A=0$ the vector transformations are the scalar ones with $N=d$). $\dm_K$, $E_K$, $\dm_k$, $E_k$ are ``the same expressions'' \cite[(1.26), (1.28)]{Ba3} with $G_K$, resp.\ $G_k$, in place of $C_0$ and $C$: the \cite[(1.7)]{Ba3} vertices inside them keep their numerical coefficients $\dm_1$, $\dm_{(4)}$, and only the lines are replaced (\cref{conv:a06d:LOC}; in \cite[(1.23)]{Ba3} $\dm_1$ is a symbol for a number). Write $\dm[G,G^v]$, $E[G,G^v]$ for the class sums \cite[(1.23)--(1.24)]{Ba3} with the lines $(G,G^v)$, so that $\dm=\dm[C_0,C]$, $\dm_K=\dm[G_K,G^v_K]$, $\dm_{\mathrm{fin}}=\dm-\dm_K$, $\dm_{K,k}=\dm_K-\dm_k$, and likewise for $E$. Superscripts $(2)$, $(4)$ denote the parts of order $\alpha+2\beta=2$, $4$ (the orders $(3,0)$, $(1,1)$ vanish, $\alpha$ odd).

\begin{lemma}[operator monotonicity]\label{lem:a06e:B0}
For $1\le k\le K$ and every $x$: \textup{(a)} $C_{M_k}\le G_k\le C_0$ as operators, $C_M\defeq(-\Delta+m^2+M^2)^{-1}$, $M_k^2\defeq al_k^{-2}$; \textup{(b)} $0\le R_k(x,x)\le c_Rl_k^{-1}$ and $|G_{K,k}(x,x)|\le c_Rl_k^{-1}$, $c_R\defeq C(m)a^{1/2}$; \textup{(c)} the same for the vector kernels.
\end{lemma}

\begin{proof}
(a) $0\le a_kl_k^{-2}P_k\le M_k^2$ and inversion is operator-monotone decreasing. Diagonal entries are monotone ($K(x,x)=\ip{e_x}{Ke_x}$, $e_x\defeq\eps^{-3}\mathbf 1_x$), so $C_{M_k}(0)\le G_k(x,x)\le C_0(0)$ and $0\le R_k(x,x)\le C_0(0)-C_{M_k}(0)=\ell_T^{-3}\sum_pM_k^2\hat C_0(p)(\hat\Delta(p)+m^2+M_k^2)^{-1}\le CM_k^2\int_{\R^3}((p^2+m^2)(p^2+M_k^2))^{-1}dp\le C'M_k$ (\cref{lem:a06e:M1}). (b) $G_{K,k}(x,x)=G_K(x,x)-G_k(x,x)$ lies between $C_{M_K}(0)-C_0(0)\ge-c_Rl_K^{-1}$ and $C_0(0)-C_{M_k}(0)$, and $l_K\ge l_k$.
\end{proof}

\begin{lemma}[the slice ladder]\label{lem:a06e:SL}
Let $\sigma,\sigma'\in\{0,1\}$ count forward differences in the first and second variable, $y\defeq x-x'$ (torus distance). \textup{(a)} $G_k=\sum_{j<k}S_j$ and $R_k=R_K+\sum_{j=k}^{K-1}S_j$. \textup{(b)} For $0\le j<K$, $|\partial^\sigma\partial'^{\sigma'}S_j(x,x')|\le C_Kl_j^{-1-\sigma-\sigma'}e^{-\delta_K|y|/l_j}$, with $C_K$, $\delta_K$ free of $j$, $\eps$, $K$ and the torus. \textup{(c)} $|\partial^\sigma\partial'^{\sigma'}G_k(x,x')|\le C(\eps+|y|)^{-1-\sigma-\sigma'}e^{-\delta_K|y|/(2l_k)}$, and $|\partial^\sigma\partial'^{\sigma'}C_0(x,x')|\le C(\eps+|y|)^{-1-\sigma-\sigma'}e^{-2c_*|y|}$ with $c_*=c_*(m)>0$ (several differences in one variable allowed, $\sigma+\sigma'\le3$), on every torus with sides $\ge1$. \textup{(d)} The same for the vector kernels.
\end{lemma}

\begin{proof}
(a) is telescoping. (b) For $j\ge1$ the slice identity \cite[(1.27)]{Ba3}: $S_j=a_j^2l_j^{-4}G_jQ_j^*C^{(j)}Q_jG_j$, with $Q_j$ the block average and $C^{(j)}=((\Delta^{(j)}+aL^{-2}P)|_{T^{(j)}})^{-1}$ the step-$j$ fluctuation covariance. In scale-$j$ units the legs $G_jQ_j^*\delta_{\hat y}$, their differences and H\"older quotients (over separations $\le2$) obey $c_0e^{-\delta_0\dist(\hat x,B_j(\hat y))}$, with $c_1(\alpha)e^{2\delta_0}$ in place of $c_0$ for the H\"older quotients of exponent $\alpha$ (the printed bound decays in $\dist(\{\hat x,\hat x'\},\hat y)\ge\dist(\hat x,B_j(\hat y))-2$), by \cite[Lemma~2.4, (2.35)--(2.36)]{Ba4} on the periodic box $T^{(j)}$ (a union of unit blocks, the torus being aligned with the blocks), in the form of \cref{thm:a06:A5}: the constants are free of $j$ and of the torus, and no depth condition arises (this is \cite[Prop.~2.1, (2.24)--(2.25)]{Ba1} at $A=0$ with the source $Q_j^*\delta_{\hat y}$, $\norm{Q_j^*\delta_{\hat y}}_\infty=1$); and $|C^{(j)}(\hat y,\hat y')|\le c_2e^{-\theta_0|\hat y-\hat y'|}$ by \cite[(2.37)]{Ba4} with the $j$-free constants of \cref{thm:a06:A5}, i.e.\ \cref{lem:a06A:blockinv}(c) at $A=0$ with $\Lambda=T^{(j)}$, whose proof uses, of the geometry in (H1), only that $T^{(j)}$ is a union of $L$-blocks (the blocks of scale $j+1\le K$); keep half of each exponent for the two block sums (\cref{lem:a06d:V11}) and use $|\hat x-\hat x'|\le\dist(\hat x,B(\hat y))+|\hat y-\hat y'|+\dist(\hat x',B(\hat y'))+2$; rescaling to $\eps$-units gives $l_j^{-1}$ for the kernel and $l_j^{-1}$ per difference, and $\delta_K\defeq\min(\delta_0,\theta_0)/2$. For $j=0$, $S_0=G_1$ in scale-$1$ units has the fixed spacing $1/L$: the $\ell^2$ decay of $G_1$ at rate $\theta_0$ (the Combes--Thomas bound of group (A), \cref{lem:a06A:CT} for $T=G_1^{-1}\ge\gamma_\emptyset$ by \cref{lem:a06A:coer}, at the fixed spacing $1/L$ and with no depth restriction) gives the pointwise bounds with a constant depending on $L$, and the rescaling is as before. (c) Sum (b) over $j<k$: the slices with $l_j\ge|y|$ form a geometric sum $\le C\max(\eps,|y|)^{-1-\sigma-\sigma'}$, and those with $l_j<|y|$ give $l_j^{-a}e^{-\delta_K|y|/l_j}\le C_a|y|^{-a}(l_j/|y|)$, again geometric; split $e^{-\delta_K|y|/l_j}\le e^{-\delta_K|y|/(2l_j)}e^{-\delta_K|y|/(2l_k)}$ for $j<k$. For $C_0$: on $\eps\Z^3$ shift the contour of the Fourier integral of $(\hat\Delta(p)+m^2)^{-1}$ by $i\tau$, $\tau\defeq m/4$, in the direction of the largest coordinate of $y$ (for $\eps\tau\le1$, $\Re\hat\Delta(p+i\tau e)\ge\hat\Delta(p)-1.1\tau^2$), together with the short-distance bound $\int_{B_\eps}|p|^{\sigma+\sigma'}(|p|^2+m^2)^{-1}dp\le C\eps^{-1-\sigma-\sigma'}$; the torus kernel is the periodisation, and the images at distance $\ge\ell_T/2\ge\frac12$ are summable by the mass. (d) is the same with $\lmu$.
\end{proof}

\begin{lemma}[the remainder]\label{lem:a06e:RK}
On every torus aligned with the blocks with all sides $\ge1$, for $1\le a\le K$ and $\sigma,\sigma'\in\{0,1\}$: $|\partial^\sigma\partial'^{\sigma'}R_a(x,x')|\le C_{R}\max(l_a,|y|)^{-1-\sigma-\sigma'}e^{-c_*|y|}$, $C_R=C(m,L,a,C_K,\delta_K)$, free of $\eps$, $a$, $K$ and the torus; the same for $R^v_a$.
\end{lemma}

\begin{proof}
Case $|y|\ge l_a$: $R_a=C_0-G_a$, and \cref{lem:a06e:SL}(c) bounds both kernels by $C(\eps+|y|)^{-1-\sigma-\sigma'}$ times $e^{-2c_*|y|}$, resp.\ $e^{-\delta_K|y|/(2l_a)}\le e^{-\delta_K|y|/2}$ ($l_a\le1$); take $c_*\le\delta_K/4$. Case $|y|<l_a$: $R_a=G_a(a_al_a^{-2}P_a)C_0$ \cite[(1.25)]{Ba3}, whose kernel is $a_al_a^{-5}\sum_B(\int_B\partial^\sigma G_a(x,\cdot))(\int_BC_0\partial'^{\sigma'*}(\cdot,x'))$, the sum over the blocks $B$ of side $l_a$; by \cite[(2.35)]{Ba4} at scale $a$ in the form of \cref{thm:a06:A5} ($\mathbf 1_B=Q_a^*\delta_{\hat y}$ in scale-$a$ units; the torus is a periodic box that is a union of the blocks of side $l_a$), $|\int_B\partial^\sigma G_a(x,\cdot)|=|(\partial^\sigma G_a\mathbf 1_B)(x)|\le Cl_a^{2-\sigma}e^{-\delta_0\dist(x,B)/(4l_a)}$, by \cref{lem:a06e:SL}(c) $\int_B|C_0\partial'^{\sigma'*}(\cdot,x')|\le Cl_a^{2-\sigma'}$, and $\sum_Be^{-\delta_0\dist(x,B)/(4l_a)}\le C$; so the kernel is $\le Cl_a^{-1-\sigma-\sigma'}$, and $e^{-c_*|y|}\ge e^{-c_*}$.
\end{proof}

\begin{proposition}[the order-$2$ infrared remainders]\label{prop:a06e:DM1}
For $a,b\in\{1,\dots,K\}$ with $|a-b|\le1$ let $\Delta_{a,b}(x)\defeq\dm_1-\dm_1[G_a,G^v_b](x)$. Then $|\Delta_{a,b}(x)|\le C_Dl_{\min(a,b)}^{-1}$, $C_D=C(N,m,\lmu,e,\lambda)$, uniformly in $x$, $\eps$ and the torus. In particular ${\dm}^{(2)}_{\mathrm{fin}}=\Delta_{K,K}=O(l_K^{-1})$ and ${\dm}^{(2)}_{K,k}=\Delta_{k,k}-\Delta_{K,K}=O(l_k^{-1})$.
\end{proposition}

\begin{proof}
By \cite[(1.23)]{Ba3}, $\dm_1[G,G^v](x)=-4(N+2)\lambda G(x,x)+e^2dq^2G^v(x,x)+b(G,G^v)(x)$ with $b(G,G^v)(x)\defeq-e^2q^2\sum_{x'}\eps^3\sum_\mu(\partial_\mu G\partial'^*_\mu)(x,x')G^v(x,x')$. (1) Tadpoles: $-4(N+2)\lambda R_a(x,x)$ and $e^2dq^2R^v_b(x,x)$ are bounded by \cref{lem:a06e:B0}. (2) $b(C_0,C)-b(G_a,G^v_b)=b(R_a,C)+b(G_a,R^v_b)$ by bilinearity. (3) $b(R_a,C)$: by \cref{lem:a06e:RK} $|\partial\partial'R_a|\le C\max(l_a,|y|)^{-3}$ and $|C(x,x')|\le C(\eps+|y|)^{-1}$ (\cref{lem:a06e:SL}(c), no additive constant), so $\sum_{x'}\eps^3|\cdot|\le C[l_a^{-3}\int_{|y|\le l_a}(\eps+|y|)^{-1}dy+\int_{|y|\ge l_a}|y|^{-4}dy]\le Cl_a^{-1}$. (4) $b(G_a,R^v_b)$, zero sum: $\sum_{x'}\partial'^*_{x'}F(x')=0$ on the torus, so $R^v_b(x,x')$ may be replaced by $R^v_b(x,x')-R^v_b(x,x)$; integrating $|\partial'R^v_b|\le C\max(l_b,s)^{-2}$ along a lattice path gives $|R^v_b(x,x')-R^v_b(x,x)|\le C\min(|y|l_b^{-2},l_b^{-1})$. For $|y|\le l_b$, $|\partial\partial'G_a|\le C(\eps+|y|)^{-3}$ gives $Cl_b^{-2}\int_{|y|\le l_b}|y|^{-2}dy\le Cl_b^{-1}$; for $|y|>l_b$, $|\partial\partial'G_a|\le C(\eps+|y|)^{-3}e^{-\delta_K|y|/(2l_a)}$ gives $Cl_b^{-1}\int_{|y|>l_b}|y|^{-3}e^{-\delta_K|y|/(2l_a)}dy\le Cl_b^{-1}(1+\ln_+(l_a/l_b))\le C(1+\ln L)l_{\min(a,b)}^{-1}$. (Without the zero-sum step the bound would carry $\ln(l_b/\eps)$.)
\end{proof}

\subsubsection*{The multi-cut expansion and the cut-line mechanism}

Let $F[G]$ denote a class sum of \cite[(1.23)]{Ba3} (two-point, rooted at $x$) or \cite[(1.24)]{Ba3} (vacuum) with all lines in a kernel pair $G=(G,G^v)$, written as Proposition~1 writes it: the sum over the forest families $\{G_{\mathcal F}\}$ of its graphs, $\mathcal F$ a forest of divergent subgraphs, each maximal member contracted to its counterterm vertex, where a subgraph is a set of lines with the vertices it touches and its degree \cite[(2.1)]{Ba3} counts only its own lines, and the coefficients of the counterterm vertices are fixed numbers (\cref{conv:a06d:LOC}).

\begin{lemma}[the multi-cut expansion]\label{lem:a06e:MC}
Let $G=H+D$ be kernel pairs. \textup{(a)} $F[H+D]-F[H]=\sum_\Gamma\sum_{\emptyset\ne\Lambda\subset\mathrm{lines}(\Gamma)}(\Gamma\setminus\Lambda)_{\mathrm{ren}}[H;D\text{ on }\Lambda]$, where $(\Gamma\setminus\Lambda)_{\mathrm{ren}}\defeq\{G_{\mathcal F}:\mathcal F$ a forest of divergent subgraphs of $\Gamma$ none of which contains a line of $\Lambda\}$, every line not in $\Lambda$ carries $H$, and every line of $\Lambda$ is cut and carries $D$ as a two-variable external field. \textup{(b)} $(\Gamma\setminus\Lambda)_{\mathrm{ren}}$ is exactly the forest family of the cut graph $\Gamma\setminus\Lambda$, and it factorises over the connected components (clusters) of $\Gamma\setminus\Lambda$. \textup{(c)} With $H=G_{j+1}$, $D=-S_j$: $m_j\defeq F[G_{j+1}]-F[G_j]=\sum_\Gamma\sum_{\Lambda\ne\emptyset}(-1)^{|\Lambda|+1}(\Gamma\setminus\Lambda)_{\mathrm{ren}}[G_{j+1};S_j\text{ on }\Lambda]$; every uncut line carries $G_{j+1}$.
\end{lemma}

\begin{proof}
In $G_{\mathcal F}$ the lines inside members of $\mathcal F$ are contracted, so the remaining lines are those of $\Gamma\setminus\bigcup\mathcal F$; expanding each as $H+D$ gives $\sum_{\Lambda\subset\mathrm{lines}(G_{\mathcal F})}G_{\mathcal F}[H\text{ off }\Lambda,D\text{ on }\Lambda]$, the term $\Lambda=\emptyset$ being $G_{\mathcal F}[H]$; counterterm vertices are vertices, not lines, and are not expanded. For fixed $\Lambda$, the forests that occur are those disjoint from $\Lambda$. (b) For a subgraph containing no line of $\Lambda$ the degree counts only its own lines, so it is the same in $\Gamma$ and in $\Gamma\setminus\Lambda$; hence the divergent subgraphs of $\Gamma\setminus\Lambda$ are exactly those of $\Gamma$ avoiding $\Lambda$. A divergent subgraph is connected, so it lies in one cluster; forests and values factorise over clusters, which are joined only through the $D$-fields on $\Lambda$. (c) is (a) with $G_j=G_{j+1}-S_j$.
\end{proof}

Proposition~1 is multilinear in the coefficients of the \cite[(1.7)]{Ba3} vertices: a \cite[(1.7)]{Ba3} vertex of order $\mathrm{ord}(v)=\alpha+2\beta$ whose coefficient is a fixed number $c_v$ (in the units of the step, couplings included) enters the bound \cite[(1.33)]{Ba3} through $\sup|c_v|\,t^{-\mathrm{ord}(v)/2}$ in place of the factor $t^{\mathrm{ord}(v)/2}$ that the coupling count $e_k^{d_v}\lambda_k^{d_s}$ assigns to it (\cite[p.~421]{Ba3}: ``the vertices (1.7) give also a positive power of $L^k\eps$ depending on a term in (1.29)''). So a localised family of order $o$ whose \cite[(1.7)]{Ba3} vertices have coefficients $c_v$ is bounded by $O(1)t^{o/2}\prod_v(\sup|c_v|t^{-\mathrm{ord}(v)/2})e^{-\delta_0d(\Lambda)}\prod\norm{h\Phi}_{1,\alpha_0}$. This is the form in which Proposition~1 is used below and in Steps~3(D$_1$)--3(D$_2$).

\begin{lemma}[the cut-line mechanism]\label{lem:a06e:CUT}
Fix $0\le j<K$ and a class $F$ of graphs of order $o$ of \cite[(1.23)]{Ba3} or \cite[(1.24)]{Ba3}, and let $m_j=F[G_{j+1}]-F[G_j]$. \textup{(a)} Each term of \cref{lem:a06e:MC}(c) is a product of clusters, renormalised families at step $j+1$ with all lines $G_{j+1}$, joined through the cut lines, which are two-variable external fields: in scale-$(j+1)$ units $|\partial^\sigma\partial'^{\sigma'}S_j(x,x')|\le C_KL^{1+\sigma+\sigma'}e^{-\delta_KL|x-x'|}$, and the two-variable norm $\norm{\cdot}_{1,\alpha_0}$ of $hS_jh'$ for localisations at boxes at distance $r$ is $\le C_KL^{2+\alpha_0}e^{-\delta_KLr}$ (\cref{lem:a06e:SL}(b); for $j=0$ its $L$-dependent constant). \textup{(b)} The only \cite[(1.7)]{Ba3} vertices in a two-point class of order $\le4$ or a vacuum class are $\dm_1$ and $\dm_{(4)}$ vertices; in a cluster with lines $G_{j+1}$ the non-renormalising part of a $\dm_1$ vertex is $\Delta_{j+1,j+1}=O(l_{j+1}^{-1})$ (\cref{prop:a06e:DM1}), so in scale-$(j+1)$ units its coefficient is $O(t_{j+1})$: it enters Proposition~1 with its printed power $t^{\mathrm{ord}/2}=t$. \textup{(c)} Proposition~1 per cluster, the clusters joined through the cut fields as in \cref{lem:a06d:V10}, bounds each localised term by $O(1)C_K^{|\Lambda|}L^{c|\Lambda|}t_{j+1}^{o/2-\sum_v(\mathrm{ord}(v)/2-1)^+}e^{-\delta''d(\{\square\})}$, the sum over the counterterm vertices $v$ of the class (a $\dm_{(4)}$ vertex carries $t\,q(t)$, $q(t)=(1+\ln t^{-1})^5$, i.e.\ $t^{1-\beta}$ for every $\beta>0$, not $t^2$: \cref{thm:a06e:T0}). \textup{(d)} Summing over localisations (\cref{lem:a06d:V11}), with the root cube fixed: for a two-point class without $\dm_{(4)}$ vertices, $|m_j(x)|\le Cl_j^{o/2-2}$; for a vacuum class, $|m_j|\le Ct_j^{o/2}|T^{(j)}|$; with the exponent of \textup{(c)}, times $q(t_{j+1})$ per $\dm_{(4)}$ vertex, when $\dm_{(4)}$ vertices occur. The constants are free of $\eps$, $\eps_0$, $j$, $K$ and the torus.
\end{lemma}

\begin{proof}
(a) is \cref{lem:a06e:MC}(c) and the rescaling of \cref{lem:a06e:SL}(b) to scale $j+1$, where $S_j$ has side $l_j/l_{j+1}=1/L$. (b) In a two-point class of order $\le4$ the \cite[(1.7)]{Ba3} vertices have order $\le2$ (a $\dm_{(4)}$ vertex alone is the whole order); in a cluster with lines $G_{j+1}$, every $\dm_1$ vertex is paired with the order-$2$ two-point subgraphs of the same class with lines $G_{j+1}$, and the pair sums to $\dm_1-\dm_1[G_{j+1},G^v_{j+1}]=\Delta_{j+1,j+1}$ at that site, plus renormalised two-point insertions, which Proposition~1 bounds. (c) Proposition~1 with the external fields of (a), in the multilinear form above; the cut fields tie the clusters, and the tree of the whole term has length at most the sum of the cluster trees and the cut-line distances. (d) In scale-$(j+1)$ units the two-point coefficient is $l_{j+1}^2m_j(x)$, bounded by $Ct_{j+1}^{o/2}$ after the sum over localisations with the root fixed; for a vacuum class the root ranges over the $|T^{(j+1)}|$ boxes. The factor $t_{j+1}^{(\cdot)}=L^{(\cdot)}t_j^{(\cdot)}$ changes constants only; the number of terms is $\le2^{I}$, $I$ the number of lines.
\end{proof}

\begin{theorem}[{$\dm_{K,k}=O((L^k\eps)^{-1})$, \cite[(1.28)]{Ba3}}]\label{thm:a06e:DM}
For $0\le k<K$ and every $x$: \textup{(a)} $\dm_{K,k}=\sum_{j=k}^{K-1}m_j$ with $m_j\defeq\dm_{j+1}-\dm_j$, the graphs of \cite[(1.23)]{Ba3} with all lines in $\{G_j,S_j\}$ and at least one $S_j$ line. \textup{(b)} Order $2$: for $j\ge1$, $m^{(2)}_j=\Delta_{j,j}-\Delta_{j+1,j+1}$, so $|m^{(2)}_j|\le2C_Dl_j^{-1}$; for $j=0$, $m^{(2)}_0=\dm_1-\Delta_{1,1}$ and $|\dm_1|\le c_ml_0^{-1}$ (\cref{thm:a06e:CT1}). \textup{(c)} Order $4$: $|m^{(4)}_j(x)|\le C_4$ (\cref{lem:a06e:CUT} with $o=4$, $l_j^{o/2-2}=1$; the only \cite[(1.7)]{Ba3} vertices are $\dm_1$ vertices). \textup{(d)} Hence $|\dm_{K,k}(x)|\le\sum_{j\ge k}(2C_Dl_j^{-1}+C_4)\le Cl_k^{-1}+C_4(K-k)\le C'l_k^{-1}$, since $K-k\le\log_L(1/l_k)\le l_k^{-1}$ ($l_K\le\eps_0\le1$); so $\dm_{K,k}(x)(L^k\eps)^2=O(L^k\eps)$, exactly as \cite[p.~418]{Ba3} asserts, with a constant free of $\eps$, $\eps_0$, $k$, $K$ and the torus. \textup{(e)} ${\dm}^{(2)}_{\mathrm{fin}}=\Delta_{K,K}=O(l_K^{-1})$. The order-$(0,2)$ part of $\dm_{K,k}$ obeys $|\cdot|\le C\lambda^2(1+\ln l_k^{-1})$ without Proposition~1.
\end{theorem}

\begin{proof}
(a) is telescoping under \cref{conv:a06d:LOC}; (b) and (e) are \cref{prop:a06e:DM1}; (c) is \cref{lem:a06e:CUT}(d); (d) is the sum. For the last clause: the order-$\lambda^2$ graphs are the sunset, the tadpole on a tadpole and the tadpole with a $\dm_{(0,1)}$ vertex, the last two forming the class ``$\lambda$-tadpole whose loop line carries the blob $4(N+2)\lambda(C_0(0)-G(x',x'))$''; the sunset part is $16(2N+4)\lambda^2\sum_{x'}\eps^3[G_K^3-G_k^3](x,x')$, bounded by \cref{lem:a06e:SL} through $|G_{K,k}(x,x')|\le\sum_{j=k}^{K-1}Cl_j^{-1}e^{-\delta_K|y|/l_j}$ and $|G_K|,|G_k|\le C(\eps+|y|)^{-1}$ by $\sum_{j=k}^{K-1}Cl_j^{-1}\int(\eps+|y|)^{-2}e^{-\delta_K|y|/l_j}dy\le C(K-k)$; the tadpole class is $C_N\lambda^2\sum_{x'}\eps^3[(G_K^2-G_k^2)R_K-G_k^2G_{K,k}(x',x')]$, with $|R_K|\le c_Rl_K^{-1}$, $|G_{K,k}(x',x')|\le c_Rl_k^{-1}$ (\cref{lem:a06e:B0}), bounded by $Cl_K^{-1}\sum_{j<K}l_j^{-1}l_j^2+Cl_k^{-1}l_k\le C$.
\end{proof}
